\chapter{Teoria Cherna--Weila i dowód wzoru Cherna--Gaussa--Bonneta}
\label{ch:chern-weil}
\index{teoria Cherna--Weila@Teoria Cherna--Weila}

W Sekcji~\ref{sec:gauss-bonnet} całka z Pfaffianu
krzywizny została sformułowana jako wynik
zewnętrzny. Pokażemy, dlaczego krzywizna
reprezentuje klasę Eulera i skąd bierze się
poprawka na brzegu. Ten sam mechanizm określa
klasy Pontriagina potrzebne przy sferach Milnora.

\section{Wielomiany niezmiennicze i transgresja}
\label{sec:chern-weil-26}

Niech $E\to M$ będzie zorientowaną rzeczywistą wiązką rzędu $q$
z metryką i koneksją metryczną $\nabla$.
W dodatniej ramce ortonormalnej $(e_1,\ldots,e_q)$
z Definicji~\ref{def:ramka-lokalna} ustalamy konwencję
$\nabla e_j=\sum_i\omega_{ij}e_i$.
Różniczkując $\langle e_i,e_j\rangle=\delta_{ij}$,
otrzymujemy $0=\omega_{ji}+\omega_{ij}$:
macierz jednoform $\omega$ jest skośnie symetryczna.
Dla macierzy form stosujemy iloczyn
$(\eta\wedge\xi)_{ij}=\sum_k\eta_{ik}\wedge\xi_{kj}$.
Kolejność czynników jest istotna, ponieważ mnożymy także macierze.

W tej ramce koneksja działa na kolumnę współczynników sekcji
według $\nabla v=\dd v+\omega v$.
Po ponownym zastosowaniu różniczki kowariantnej mamy
$\dd(\omega v)=\dd\omega\,v-\omega\wedge\dd v$;
drugi składnik znosi wyraz $\omega\wedge\dd v$.
Stąd krzywizna ma macierz
$\Omega=\dd\omega+\omega\wedge\omega$.
Jej wyrazy są dwusformami i spełniają
$\Omega_{ij}(X,Y)=\langle R^\nabla(X,Y)e_j,e_i\rangle$,
zgodnie z konwencją Rozdziału~\ref{ch:krzywizna}.
Jeśli zmienimy ramkę przez $\widetilde e=eQ$,
$Q:U\to\mathrm{SO}(q)$, reguła iloczynu daje
\[
 \widetilde\omega=Q^{-1}\omega Q+Q^{-1}\dd Q,\qquad
 \widetilde\Omega=Q^{-1}\Omega Q.
\]
Pierwszy wzór otrzymujemy z $\nabla(eQ)=(\nabla e)Q+e\,\dd Q$.
Drugi wynika z tensorowości krzywizny:
$R^\nabla(X,Y)(eQ)=(R^\nabla(X,Y)e)Q$.

\begin{definicja}[Wielomian niezmienniczy]
\index{wielomian niezmienniczy@Wielomian niezmienniczy}
\label{def:invariant-polynomial-26}
Symetryczne $r$-liniowe odwzorowanie
$P:\mathfrak{so}(q)^r\to\R$ nazywamy
\emph{niezmienniczym}, gdy
$P(gA_1g^{-1},\ldots,gA_rg^{-1})
=P(A_1,\ldots,A_r)$ dla $g\in\mathrm{SO}(q)$.
Zakładamy $r\geq1$. Przez $P(\Omega^r)$ rozumiemy
$P(\Omega,\ldots,\Omega)$ z $r$ argumentami,
a nie potęgę macierzy $\Omega$.
Wyjaśnijmy podstawianie form do $P$.
Jeśli w bazie macierzy $(T_a)$ zapiszemy
$\eta_j=\sum_a T_a\eta_j^a$, to kładziemy
\[
 P(\eta_1,\ldots,\eta_r)
 =\sum_{a_1,\ldots,a_r}
 P(T_{a_1},\ldots,T_{a_r})\,
 \eta_1^{a_1}\wedge\cdots\wedge\eta_r^{a_r}.
\]
W szczególności $P(\Omega^r)$ jest formą stopnia $2r$.
Niezmienniczość i wzór na $\widetilde\Omega$ pokazują,
że lokalne formy są takie same na przecięciach map;
otrzymujemy więc formę globalną.
\end{definicja}

Przykładem jest ślad $P(A,B)=\operatorname{tr}(AB)$:
symetrię daje $\operatorname{tr}(AB)=\operatorname{tr}(BA)$,
a niezmienniczość wynika z
$\operatorname{tr}(gABg^{-1})=\operatorname{tr}(AB)$.
Dla $q=2m$ polaryzacja Pfaffianu
z~\eqref{eq:gb-pfaff-def} jest
niezmienniczym wielomianem stopnia $m$.
Ustalamy jej normalizację przez
$P_{\operatorname{Pf}}(B,\ldots,B)=\operatorname{Pf}(B)$.
Równoważnie, $m!P_{\operatorname{Pf}}(B_1,\ldots,B_m)$
jest współczynnikiem przy $t_1\cdots t_m$ w
$\operatorname{Pf}(t_1B_1+\cdots+t_mB_m)$.
Taka normalizacja będzie potrzebna przy wyznaczaniu
stałej w formie kątowej.

\begin{twierdzenie}[Cherna--Weila i transgresja]
\label{thm:chern-weil-26}
Forma $P(\Omega^r)$ jest zamknięta,
a jej klasa de Rhama nie zależy od koneksji.
Dokładniej, dla dwóch koneksji metrycznych na tej samej wiązce
z ustaloną metryką i dla $\nabla^t=\nabla^0+tA$,
$A=\nabla^1-\nabla^0$, zachodzi
\begin{equation}\label{eq:CW-transgression-26}
 P(\Omega_1^r)-P(\Omega_0^r)
 =\dd\left(r\int_0^1
       P(A,\Omega_t^{r-1})\,\dd t\right).
\end{equation}
Konstrukcja jest naturalna względem cofnięcia wiązki.
\end{twierdzenie}
\begin{proof}
\emph{Krok 1: różnica koneksji i różniczka kowariantna.}
Dla funkcji $f$ i sekcji $v$ wyrazy $(\dd f)v$ w
$\nabla^1(fv)-\nabla^0(fv)$ znoszą się.
Zatem $A$ jest liniowe nad funkcjami względem $v$,
a względem kierunku różniczkowania jest jednoformą.
W ramce jego macierz to $\omega_1-\omega_0$.
Wyrazy $Q^{-1}\dd Q$ we wzorze zmiany ramki także się znoszą,
więc $\widetilde A=Q^{-1}AQ$.
Otrzymujemy globalną jednoformę o wartościach
w skośnie symetrycznych endomorfizmach $E$.
Ścieżka $\nabla^t$ jest więc globalna i metryczna.

Dla endomorfizmowej $k$-formy $\eta$ definiujemy
$D_t\eta=\dd\eta+\omega_t\wedge\eta-(-1)^k\eta\wedge\omega_t$.
Dla $k=0$ jest to zwykły komutator z macierzą koneksji;
znak w ogólnym wzorze pochodzi z przeniesienia jednoformy
przez $k$ różniczek. Rozwinięcie krzywizny daje
\begin{align*}
 \Omega_t&=\Omega_0+tD_0A+t^2A\wedge A,\\
 \frac{\dd}{\dd t}\Omega_t
 &=\dd A+\omega_t\wedge A+A\wedge\omega_t=D_tA.
\end{align*}
Tożsamość Bianchiego sprawdzamy wprost:
\[
 \dd\Omega_t
 =\dd\omega_t\wedge\omega_t-\omega_t\wedge\dd\omega_t,
 \qquad
 D_t\Omega_t
 =\dd\Omega_t+\omega_t\wedge\Omega_t-\Omega_t\wedge\omega_t=0.
\]
Po podstawieniu $\Omega_t=\dd\omega_t+\omega_t\wedge\omega_t$
znoszą się zarówno wyrazy z $\dd\omega_t$,
jak i dwa wyrazy z trzema czynnikami $\omega_t$.

\emph{Krok 2: dlaczego znikają komutatory.}
Podstawmy $g=\exp(sB)$ do niezmienniczości $P$
i zróżniczkujmy po $s$ w zerze. Otrzymujemy
\[
 \sum_{j=1}^r
 P(B_1,\ldots,[B,B_j],\ldots,B_r)=0.
\]
Dla form $\eta_j$ stopni $k_j$ zwykła reguła iloczynu daje
\[
 \dd P(\eta_1,\ldots,\eta_r)
 =\sum_{j=1}^r(-1)^{k_1+\cdots+k_{j-1}}
 P(\eta_1,\ldots,D_t\eta_j,\ldots,\eta_r).
\]
Uzasadnijmy zastąpienie $\dd\eta_j$ przez $D_t\eta_j$.
Rozpisujemy $\omega_t$ jako sumę macierzy razy jednoformy.
Po przeniesieniu takiej jednoformy na początek każdego
iloczynu pojawia się znak $(-1)^{k_1+\cdots+k_{j-1}}$,
który znosi znak stojący przed tym składnikiem.
Pozostała suma macierzowych komutatorów jest zerowa
na mocy właśnie wyprowadzonej tożsamości.

\emph{Krok 3: zamkniętość i transgresja.}
Wszystkie argumenty $\Omega_t$ mają parzysty stopień,
więc ich zamiana nie zmienia znaku iloczynu form.
Z kroku 2 i $D_t\Omega_t=0$ wynika
$\dd P(\Omega_t^r)=rP(D_t\Omega_t,\Omega_t^{r-1})=0$.
Gdy pierwszym argumentem jest $A$, dostajemy
\[
 \dd P(A,\Omega_t^{r-1})=P(D_tA,\Omega_t^{r-1}),\qquad
 \frac{\dd}{\dd t}P(\Omega_t^r)
 =rP(D_tA,\Omega_t^{r-1}).
\]
W pierwszej równości pozostałe składniki zawierają
$D_t\Omega_t=0$; w drugiej różniczkujemy kolejno $r$ argumentów,
otrzymując $r$ równych składników.
Całkujemy ostatnią równość po $t\in[0,1]$
i przenosimy $\dd$ przed całkę.
Jest to dozwolone, ponieważ lokalne współczynniki form
są gładkie w $(x,t)$, a przedział całkowania jest zwarty.
Otrzymujemy~\eqref{eq:CW-transgression-26};
różnica form jest dokładna, więc ich klasy de Rhama są równe.
Wreszcie dla gładkiego $F:N\to M$ cofnięcie komutuje
z $\dd$ i iloczynem zewnętrznym:
$F^*\Omega_\nabla=\Omega_{F^*\nabla}$ oraz
$F^*P(\Omega_\nabla^r)=P(\Omega_{F^*\nabla}^r)$.
To dowodzi naturalności.
\end{proof}

Forma w nawiasie w~\eqref{eq:CW-transgression-26}
to \emph{forma transgresji} Cherna--Simonsa:
zapisuje zmianę reprezentanta różniczkowego
przy stałej klasie kohomologii.

\begin{przyklad}[Euler i pierwsza klasa Pontriagina]
\label{prz:euler-p1-CW-26}
Dla orientowanej wiązki rzędu $2m$ określamy
$e_\nabla(E)=\operatorname{Pf}(\Omega)/(2\pi)^m$.
Zidentyfikujemy ją z topologiczną klasą
Eulera poniżej. Dla rzeczywistej wiązki
metrycznej normalizacja pierwszej formy
Pontriagina jest
\[
 p_{1,\nabla}(E)=
 -\frac{1}{8\pi^2}\operatorname{tr}
                (\Omega\wedge\Omega).
\]
Jeśli $E$ rozszczepia się na ortogonalną
sumę orientowanych $2$-płaszczyzn,
a koneksja zachowuje ten rozkład
i ma bloki krzywizny
$\Omega_i=\begin{psmallmatrix}0&\kappa_i\\-\kappa_i&0\end{psmallmatrix}$,
gdzie $\kappa_i$ są dwusformami, to
$\Omega_i\wedge\Omega_i
=\operatorname{diag}(-\kappa_i\wedge\kappa_i,
                     -\kappa_i\wedge\kappa_i)$.
Zatem $\operatorname{tr}(\Omega\wedge\Omega)
=-2\sum_i\kappa_i\wedge\kappa_i$ i
$\operatorname{Pf}(\Omega_i)=\kappa_i$.
Po podstawieniu dostajemy
$p_{1,\nabla}=\sum_i(\kappa_i/2\pi)\wedge(\kappa_i/2\pi)
=\sum_i e_\nabla(E_i)\wedge e_\nabla(E_i)$.
Kwadrat dwusformy oznacza tu jej iloczyn zewnętrzny
z nią samą; nie musi być zerowy, bo jej stopień jest parzysty.
Ogólną identyfikację można uzyskać bez
założenia takiego rozszczepienia. Z definicji
topologicznej w~\ref{sec:klasy-char-16}
klasa $p_1(E)$ jest równa
$-c_2(E\otimes_\R\C)$. Koneksja zespolona
ma tę samą macierz krzywizny $\Omega$;
ponieważ $\operatorname{tr}\Omega=0$,
składnik stopnia cztery wyrażenia
$\det(I+i\Omega/(2\pi))$ wynosi
$\operatorname{tr}(\Omega\wedge\Omega)/(8\pi^2)$.
Wyprowadźmy współczynnik. Dla macierzy $B$ o wyrazach
w przemiennym pierścieniu współczynnik przy $t^2$
w $\det(I+tB)$ jest sumą głównych minorów rzędu dwa:
\[
 \sum_{i<j}(B_{ii}B_{jj}-B_{ij}B_{ji})
 =\tfrac12\bigl((\operatorname{tr}B)^2-\operatorname{tr}(B^2)\bigr).
\]
Pierwsza suma powstaje przez wybranie dwóch czynników $tB$
w rozwinięciu wyznacznika. Po rozwinięciu prawej strony
wyrazy $B_{ii}^2$ znoszą się, a każda para $i<j$
jest policzona dwa razy przed podzieleniem przez $2$.
Dwusformy mnożą się przemiennie, więc podstawiamy
$B=i\Omega/(2\pi)$, otrzymując
$c_{2,\nabla}=(\operatorname{tr}(\Omega\wedge\Omega)
-\operatorname{tr}\Omega\wedge\operatorname{tr}\Omega)/(8\pi^2)$.
Zerowy ślad macierzy skośnie symetrycznej usuwa drugi wyraz.
Zgodność tak otrzymanej formy Cherna
z rzeczywistym obrazem topologicznej klasy
$c_2$ jest tu jawnym wejściem zewnętrznym
twierdzenia Cherna--Weila dla wiązek zespolonych.
Definicję topologiczną i kontrolny wzór
$p_1(E_i)=e(E_i)^2$ podaje
\href{https://pi.math.cornell.edu/~hatcher/VBKT/VB.pdf}
{A.\ Hatcher, \emph{Vector Bundles and K-Theory},
§3.1--3.2}; zgodność z formą krzywizny
oraz obliczenie $c_2$ opisuje
\href{https://webusers.imj-prg.fr/~olivier.biquard/m2dg2022.pdf}
{O.\ Biquarda, \emph{Differential and Riemannian Geometry},
§16.c, wzory (III.39) i (III.41)}.
\end{przyklad}

\begin{przyklad}[Niezerowa forma krzywizny przy zerowej klasie]
\label{prz:CW-plaska-wiazka-17}
Na trywialnej zorientowanej wiązce $E=\R^2\times\R^2\to\R^2$
z metryką euklidesową porównajmy koneksje
$\nabla^0=\dd$ i $\nabla^1=\dd+\omega_1$, gdzie
$\omega_1=\begin{psmallmatrix}0&x\,\dd y\\-x\,\dd y&0\end{psmallmatrix}$.
Macierz jest skośnie symetryczna, więc koneksja jest metryczna.
Każdy wyraz $\omega_1\wedge\omega_1$ zawiera
$\dd y\wedge\dd y=0$, natomiast $\dd(x\,\dd y)=\dd x\wedge\dd y$.
Zatem $\Omega_1=
\begin{psmallmatrix}0&\dd x\wedge\dd y\\-\dd x\wedge\dd y&0\end{psmallmatrix}$
i $e_{\nabla^1}(E)=\dd x\wedge\dd y/(2\pi)$.
Dla $\nabla^0$ krzywizna i forma Eulera są zerowe.
Przy $m=1$ Pfaffian wybiera wyraz o indeksach $12$,
więc forma transgresji jest równa $x\,\dd y/(2\pi)$.
Jej różniczka daje dokładnie różnicę obu form Eulera.
Forma krzywizny nie znika, ale jej klasa jest zerowa;
zgadza się to z istnieniem globalnej niezerowej sekcji $E$.

Sprawdźmy to na prostokącie $Q=[0,a]\times[0,b]$, $a,b>0$,
z orientacją $\dd x\wedge\dd y$ i dodatnio, przeciwnie do
ruchu wskazówek zegara, zorientowanym brzegiem.
Po dwóch poziomych bokach $\dd y=0$, a po lewym $x=0$.
Na prawym boku $x=a$ i $y$ rośnie od $0$ do $b$.
Dlatego
\[
 \int_Q e_{\nabla^1}
 =\frac1{2\pi}\int_0^a\int_0^b\dd y\,\dd x
 =\frac{ab}{2\pi}
 =\frac1{2\pi}\int_0^b a\,\dd y
 =\int_{\partial Q}\frac{x\,\dd y}{2\pi}.
\]
To konkretna postać transgresji i twierdzenia Stokesa.
Krzywizna w tym przykładzie dotyczy wybranej koneksji
na $E$; metryka euklidesowa podstawy służy do narysowania obszaru.
\end{przyklad}

\begin{figure}[htbp]
\centering
\begin{tikzpicture}[>=Stealth]
 \fill[manifoldblue!10] (0,0) rectangle (3.2,2.1);
 \draw[->,thick,manifoldblue] (0,0)--(3.2,0);
 \draw[->,very thick,manifoldred] (3.2,0)--(3.2,2.1);
 \draw[->,thick,manifoldblue] (3.2,2.1)--(0,2.1);
 \draw[->,thick,manifoldblue] (0,2.1)--(0,0);
 \node[below left=3pt] at (0,0) {$0$};
 \node[below=4pt] at (1.6,0) {$a$};
 \node[left=4pt] at (0,1.05) {$b$};
 \node at (1.6,1.05) {$\displaystyle\int_Q\dd x\wedge\dd y=ab$};
 \node[anchor=west,text width=5.9cm,align=left,font=\small] at (4,1.05)
  {Na prawym boku $\int_0^b a\,\dd y=ab$.\\[5pt]
   Na pozostałych bokach forma $x\,\dd y$ daje zero.
   Całka we wnętrzu przechodzi w całkę po dodatnim brzegu.};
\end{tikzpicture}
\caption{Forma dokładna może mieć niezerową całkę po obszarze
z brzegiem. W Przykładzie~\ref{prz:CW-plaska-wiazka-17}
obie strony dzielimy przez $2\pi$.
Strzałki podają orientację całego $\partial Q$.}
\label{fig:CW-transgresja-prostokat}
\end{figure}

\section{Znormalizowana forma kątowa}
\label{sec:angular-form-26}

Niech teraz $E$ ma rangę $2m$, gdzie $m\geq1$.
Na wiązce sfer jednostkowych
$\pi:S(E)\to M$ istnieje
tautologiczny przekrój $s$ wiązki
$\pi^*E$: w punkcie $(x,u)$
jest nim wektor $u\in E_x=(\pi^*E)_{(x,u)}$.
Nie wybieramy tu jednego wektora dla każdego $x\in M$:
podstawą nowej wiązki jest całe $S(E)$,
a jej punkt już zawiera wybrany $u$.
Rozszczepmy
$\pi^*E=\R s\oplus s^\perp$.
Koneksja $\nabla^s$ niech czyni
$s$ równoległym i na $s^\perp$
niech będzie rzutem $\pi^*\nabla$.
Dokładnie, jeśli $w\perp s$, to
$\nabla^s(fs+w)=\dd f\,s+\operatorname{pr}_{s^\perp}(\pi^*\nabla w)$.
Rzut zachowuje regułę Leibniza, a różniczkowanie iloczynów
skalarnych pokazuje zgodność z metryką na obu prostopadłych
składnikach. Kładąc $v=(\pi^*\nabla)s$, mamy $\langle v,s\rangle=0$,
gdyż $|s|=1$. Różnica koneksji działa według
\[
 A(X)z=\langle v(X),z\rangle s-\langle s,z\rangle v(X).
\]
Dla $z=s$ daje $A(X)s=-v(X)$, więc $\nabla^s s=0$.
Dla $z\perp s$ usuwa składową normalną pochodnej,
bo $\langle s,\pi^*\nabla_Xz\rangle=-\langle v(X),z\rangle$.
To wyprowadza podany wzór na $A=\nabla^s-\pi^*\nabla$.
Jej forma Eulera jest zerowa,
ponieważ krzywizna ma zerowy
wiersz i kolumnę w kierunku $s$.
Istotnie, $\nabla^s s=0$ pociąga $R^{\nabla^s}(X,Y)s=0$,
a skośna symetria zeruje również odpowiedni wiersz.
Każdy wyraz Pfaffianu paruje ten indeks z innym,
więc zawiera czynnik zerowy.

\begin{lemat}[Forma kątowa]
\label{lem:angular-form-26}
Niech $\Psi_\nabla$ będzie transgresją
Pfaffianu od $\pi^*\nabla$ do $\nabla^s$:
\[
 \Psi_\nabla=
 \frac{m}{(2\pi)^m}\int_0^1
 P_{\operatorname{Pf}}(A,\Omega_t^{m-1})\,\dd t,
 \qquad A=\nabla^s-\pi^*\nabla.
\]
Wtedy
$\dd\Psi_\nabla=-\pi^*e_\nabla(E)$
oraz $\int_{S(E_x)}\Psi_\nabla=1$
dla każdego włókna zorientowanego
jako brzeg dysku.
\end{lemat}
\begin{proof}
Pierwsza równość jest
\eqref{eq:CW-transgression-26},
bo forma Eulera $\nabla^s$ znika.
Obliczmy normalizację na jednym włóknie $S(E_x)\cong S^{2m-1}$.
W ortogonalnej trywializacji wyrazy macierzy
cofniętej koneksji zawierają różniczki podstawy.
Na ustalonym włóknie te różniczki są zerowe,
więc $\pi^*\nabla$ ogranicza się do płaskiej koneksji $\dd$.
Niech $u$ będzie wektorem jednostkowym.
Macierz koneksji zachowującej $u$
wynosi
$A=u\,\dd u^{\mathsf T}-\dd u\,u^{\mathsf T}$:
jest to poprzedni wzór z $v=\dd u$.
Rzeczywiście, z $u^{\mathsf T}u=1$ mamy
$\dd u^{\mathsf T}u=0$, więc $Au=-\dd u$.
W punkcie $u=e_1$ mamy
$A_{1a}=\dd u^a$, $A_{a1}=-\dd u^a$,
a $A_{ab}=0$ dla $a,b>1$.
Krzywizna $\dd+tA$ ma wtedy
część styczną
$(\Omega_t)_{ab}=(2t-t^2)
\dd u^a\wedge\dd u^b$:
$\dd A=2\,\dd u\wedge\dd u^{\mathsf T}$,
ponieważ różniczkowanie $u\,\dd u^{\mathsf T}$ daje
$\dd u\wedge\dd u^{\mathsf T}$,
a różniczkowanie $-\dd u\,u^{\mathsf T}$ daje drugi taki sam składnik.
W punkcie $e_1$, dla $a,b>1$,
$(A\wedge A)_{ab}=A_{a1}\wedge A_{1b}
=-\dd u^a\wedge\dd u^b$.
To wyjaśnia oba współczynniki, $2t$ i $-t^2$.
Ponadto $\dd u^1=0$ na przestrzeni stycznej sfery w $e_1$,
więc pierwszy wiersz i kolumna $\Omega_t$ są tam zerowe.

Policzmy dokładnie współczynnik w
$mP_{\operatorname{Pf}}(A,\Omega_t^{m-1})$.
Jest to część liniowa względem $\lambda$ w
$\operatorname{Pf}(\Omega_t+\lambda A)$;
wybieramy jeden czynnik $A$ i $m-1$ czynników $\Omega_t$,
z zachowaniem kolejności iloczynów zewnętrznych.
Indeks $1$ musi trafić do $A_{1j}$, gdyż w $\Omega_t$
pierwszy wiersz jest zerowy.
Możemy wybrać $j$ na $2m-1$ sposobów.
Pozostałe $2m-2$ indeksy parujemy na
$(2m-3)!!$ sposobów: partnera pierwszego wybieramy
na $2m-3$ sposobów, następnego na $2m-5$ i tak dalej.
Znak każdego parowania znosi znak uporządkowania
różniczek do $\dd u^2\wedge\cdots\wedge\dd u^{2m}$.
Przy wyborze $j$ czynnik $(-1)^j$ z rozwinięcia Pfaffianu
znosi się ze znakiem $(-1)^{j-2}$ przeniesienia $\dd u^j$
na jego miejsce. Wszystkie wyrazy są więc dodatnie.
Współczynnik wynosi $(2m-1)(2m-3)!!=(2m-1)!!$.
Przyjmujemy $(-1)!!=0!!=1$, aby rachunek obejmował także $m=1$.

W punkcie $e_1$ forma
$\dd u^2\wedge\cdots\wedge\dd u^{2m}$
jest dodatnią formą objętości sfery z orientacją brzegu:
zewnętrzną normalną jest $e_1$.
Obroty z $\mathrm{SO}(2m)$ przenoszą rachunek na dowolny
punkt włókna, zachowując tę orientację.
Otrzymujemy
\[
 \Psi_\nabla|_{S^{2m-1}}
 =\frac{(2m-1)!!}{(2\pi)^m}
 \left(\int_0^1(2t-t^2)^{m-1}\dd t\right)
 \operatorname{vol}_{S^{2m-1}}.
\]
Pozostają dwie stałe, które także wyprowadzimy.
Po podstawieniu $v=1-t$ oznaczmy
$I_m=\int_0^1(1-v^2)^{m-1}\,\dd v$.
Dla $m=1$ mamy $I_1=1$.
Dla $m\geq2$ całkowanie przez części, z zerowym wyrazem
$[v(1-v^2)^{m-1}]_0^1$, daje
\[
 I_m=2(m-1)\int_0^1v^2(1-v^2)^{m-2}\,\dd v
     =2(m-1)(I_{m-1}-I_m).
\]
W ostatniej równości użyliśmy $v^2=1-(1-v^2)$.
Zatem $I_m=\frac{2m-2}{2m-1}I_{m-1}$,
a indukcja daje $I_m=(2m-2)!!/(2m-1)!!$.

Objętość sfery obliczamy z całki Gaussa.
Niech $G=\int_\R e^{-x^2}\,\dd x$.
Całka jest dodatnia i skończona:
dla $|x|\geq1$ mamy $e^{-x^2}\leq e^{-|x|}$,
a na $[-1,1]$ funkcja podcałkowa jest ograniczona.
Fubini i współrzędne biegunowe na płaszczyźnie dają
\[
 G^2=\int_{\R^2}e^{-(x^2+y^2)}\,\dd x\,\dd y
 =\int_0^{2\pi}\int_0^\infty e^{-r^2}r\,\dd r\,\dd\theta
 =\pi,
\]
bo podstawienie $z=r^2$ daje
$\int_0^\infty e^{-r^2}r\,\dd r=\frac12\int_0^\infty e^{-z}\,\dd z=1/2$.
W wymiarze $2m$ analogiczna zmiana zmiennych ma czynnik
$r^{2m-1}$: promień skaluje każdy z $2m-1$ kierunków
stycznych do sfery przez $r$.
Stąd
\[
 \pi^m=G^{2m}
 =\operatorname{vol}(S^{2m-1})\int_0^\infty e^{-r^2}r^{2m-1}\,\dd r
 =\frac{\operatorname{vol}(S^{2m-1})}{2}
     \int_0^\infty e^{-z}z^{m-1}\,\dd z.
\]
Ostatnia całka jest równa $(m-1)!$:
dla $m=1$ wynosi $1$, a całkowanie przez części
zmniejsza wykładnik o jeden i daje czynnik $m-1$.
Wyrazy brzegowe znikają dzięki wykładniczemu zanikowi w nieskończoności.
Otrzymujemy $\operatorname{vol}(S^{2m-1})=2\pi^m/(m-1)!$.
Ostatecznie
\[
 \int_{S(E_x)}\Psi_\nabla
 =\frac{(2m-1)!!}{(2\pi)^m}
   \frac{(2m-2)!!}{(2m-1)!!}\,
   \frac{2\pi^m}{(m-1)!}
 =1,
\]
ponieważ $(2m-2)!!=2^{m-1}(m-1)!$.
\end{proof}

\begin{stwierdzenie}[Identyfikacja klasy Eulera]
\label{prop:euler-identification-26}
Klasa de Rhama $[e_\nabla(E)]$
jest rzeczywistym obrazem
topologicznej klasy Eulera $e(E)$.
\end{stwierdzenie}
\begin{proof}
Wyjaśnijmy użyty model względny. Dla inkluzji
$i:S(E)\hookrightarrow D(E)$ jego kochainy
stopnia $k$ są parami
$(\alpha,\beta)\in\Omega^k(D(E))\oplus
\Omega^{k-1}(S(E))$ z różniczką
$\dd(\alpha,\beta)=(\dd\alpha,i^*\alpha-\dd\beta)$.
Różniczka ma kwadrat zerowy, ponieważ jej druga składowa
po dwukrotnym zastosowaniu to
$i^*\dd\alpha-\dd(i^*\alpha-\dd\beta)=0$.
Istotnie, obcięcie form z $D(E)$ na jego brzeg
jest surjektywne: formę na $S(E)$ przedłużamy
przez kołnierz i mnożymy przez funkcję równą $1$ przy brzegu
oraz zerową poza kołnierzem.
Jeśli $(\alpha,\beta)$ jest zamknięta i
$i^*\widetilde\beta=\beta$, odejmujemy
$\dd(\widetilde\beta,0)=(\dd\widetilde\beta,\beta)$.
Otrzymujemy parę $(\alpha-\dd\widetilde\beta,0)$;
pierwsza forma jest zamknięta i znika po obcięciu,
bo $i^*\alpha=\dd\beta$.
Jeśli z kolei $(\alpha,0)=\dd(\gamma,\delta)$, to
$\alpha=\dd\gamma$ i $i^*\gamma=\dd\delta$.
Przedłużamy $\delta$ do $\widetilde\delta$ i kładziemy
$\gamma'=\gamma-\dd\widetilde\delta$.
Wtedy $\dd\gamma'=\alpha$ oraz $i^*\gamma'=0$.
Kompleks par jest przez to równoważny kompleksowi
form na $D(E)$ znikających po obcięciu do $S(E)$.
Z twierdzenia de Rhama i ciągów dokładnych par
wynika, że jego kohomologia jest
$H^*(D(E),S(E);\R)$; por.
\href{https://doi.org/10.1007/978-1-4757-3951-0}
{R.\ Bott i L.\ Tu, \emph{Differential Forms in Algebraic Topology},
§7}.

Oznaczmy przez $\pi_D:D(E)\to M$ projekcję wiązki dysków;
jej ograniczeniem do $S(E)$ jest dotychczasowe $\pi$.
Para $(\pi_D^*e_\nabla(E),-\Psi_\nabla)$
jest zamknięta w tym kompleksie, ponieważ
$\dd e_\nabla(E)=0$ oraz
$\dd(-\Psi_\nabla)=\pi^*e_\nabla(E)$ na $S(E)$.
Jej ewaluacja na względnym cyklu włókna
ma postać $\int_{D(E_x)}\alpha-\int_{S(E_x)}\beta$:
znak minus wynika ze Stokesa i powyższej
różniczki na parach.
Sprawdźmy, że całka względna znika na parach dokładnych:
\[
 \int_{D(E_x)}\dd\gamma
 -\int_{S(E_x)}(i^*\gamma-\dd\delta)
 =\int_{S(E_x)}\dd\delta=0.
\]
Pierwsze dwa składniki znoszą się przez Stokesa,
a ostatnia całka jest zerowa, ponieważ sfera nie ma brzegu.
Jej całka względna po każdym
$(D^{2m},S^{2m-1})$ wynosi
$0-\int_{S^{2m-1}}(-\Psi_\nabla)=1$.
Pierwszy składnik jest zerowy: ograniczenie $\pi_D$
do włókna jest stałe, a forma Eulera ma dodatni stopień $2m$.
Jest więc znormalizowaną klasą
Thoma z Twierdzenia~\ref{thm:euler-thom-16}.
Cofnięcie jej przez zerowy przekrój
jest z jednej strony topologiczną
klasą Eulera, a z drugiej klasą
formy $e_\nabla(E)$.
Dokładniej, zapomnienie o względności wysyła parę na
$[\pi_D^*e_\nabla(E)]$, a następnie
$s_0^*\pi_D^*e_\nabla(E)=e_\nabla(E)$,
bo $\pi_D\circ s_0=\id_M$.
\end{proof}

\begin{przyklad}[Forma kątowa w randze dwa]
\label{prz:angular-rank-two-17}
W dodatniej ramce ortonormalnej zapiszmy
$\omega=\begin{psmallmatrix}0&\beta\\-\beta&0\end{psmallmatrix}$.
Ponieważ $\beta\wedge\beta=0$, forma Eulera jest równa
$e_\nabla=\dd\beta/(2\pi)$.
Jednostkowy wektor ma lokalnie postać
$u=\cos\theta\,e_1+\sin\theta\,e_2$.
Koneksja czyniąca go równoległym ma wyraz
$(\omega^s)_{12}=\dd\theta$:
rzeczywiście, $\dd u=(-\sin\theta,\cos\theta)^{\mathsf T}\dd\theta$
współrzędnościowo, a iloczyn $\omega^s u$ jest jego przeciwieństwem.
Stąd $A_{12}=\dd\theta-\beta$.
Dla $m=1$ transgresja nie zawiera $\Omega_t$,
więc całka po $t\in[0,1]$ daje
\[
 \Psi_\nabla=\frac{\dd\theta-\pi^*\beta}{2\pi},
 \qquad
 \dd\Psi_\nabla=-\pi^*\frac{\dd\beta}{2\pi}.
\]
Na jednym włóknie $\pi^*\beta=0$ i dodatni obrót ma
$\theta$ rosnące od $0$ do $2\pi$.
Dlatego $\int_{S(E_x)}\Psi_\nabla
=\frac1{2\pi}\int_0^{2\pi}\dd\theta=1$.
Kąt $\theta$ jest tylko lokalną funkcją,
ale forma jest globalna.
Istotnie, obrót ramki o $\varphi$ daje
$\widetilde\theta=\theta-\varphi$ i
$\widetilde\beta=\beta-\dd\varphi$
ze wzoru na zmianę macierzy koneksji.
Zatem $\dd\widetilde\theta-\widetilde\beta
=\dd\theta-\beta$ na wspólnej dziedzinie.
\end{przyklad}

\begin{figure}[H]
\centering
\begin{tikzpicture}[>=Stealth]
 \fill[manifoldblue!6] (0,0) circle(1.4);
 \draw[manifoldblue!60] (0,0) circle(1.4);
 \draw[->,manifoldblue!40] (-1.65,0)--(1.85,0) node[right] {$e_1$};
 \draw[->,manifoldblue!40] (0,-1.6)--(0,1.8) node[above] {$e_2$};
 \draw[->,very thick,manifoldred] (0,0)--(45:1.4);
 \fill[manifoldblue!80!black] (45:1.4) circle(2pt);
 \node[above right=3pt] at (45:1.4) {$u$};
 \draw[manifoldgold!85!black] (.55,0) arc[start angle=0,end angle=45,radius=.55];
 \node at (.77,.25) {$\theta$};
 \draw[->,thick,manifoldblue] (110:1.4) arc[start angle=110,end angle=285,radius=1.4];
 \node[anchor=west,text width=5.9cm,align=left,font=\small] at (3,0)
  {Na włóknie $\Psi=\dd\theta/(2\pi)$: pełny dodatni obrót daje całkę $1$.\\[6pt]
   Przy zmianie ramki zmieniają się zarówno $\theta$, jak i $\beta$.
   Różnica $\dd\theta-\beta$ pozostaje taka sama.};
\end{tikzpicture}
\caption{Normalizacja formy kątowej na włóknie kołowym.
Wektor $u$ jest punktem sfery jednostkowej;
strzałka na okręgu określa dodatnią orientację włókna.}
\label{fig:angular-kolo-normalizacja}
\end{figure}

\section{Dowód i poprawka brzegowa}
\label{sec:CGB-proof-26}

\begin{twierdzenie}[Cherna--Gaussa--Bonneta
z gładkim brzegiem]
\label{thm:CGB-boundary-26}
Niech $(M^{2m},g)$ będzie zwartą,
zorientowaną rozmaitością
riemannowską o gładkim brzegu,
a $\nu$ zewnętrzną jednostkową
normalną, traktowaną jako odwzorowanie
$\nu:\partial M\to S(TM)$, $x\mapsto(x,\nu_x)$.
Brzeg orientujemy regułą „zewnętrzna normalna jako pierwszy
wektor”: $(v_1,\ldots,v_{2m-1})$ jest dodatnią bazą
$T_x\partial M$, gdy $(\nu_x,v_1,\ldots,v_{2m-1})$
jest dodatnią bazą $T_xM$.
Określmy konkretną
formę brzegową
\begin{equation}\label{eq:CGB-boundary-form-26}
 \mathcal B_g:=\nu^*
       \Psi_{\nabla^{TM}}.
\end{equation}
Wtedy
\begin{equation}\label{eq:CGB-boundary-26}
 \boxed{\displaystyle
 \chi(M)=\frac{1}{(2\pi)^m}
     \int_M\operatorname{Pf}(\Omega^{TM})
       +\int_{\partial M}\mathcal B_g.}
\end{equation}
Jeśli brzeg jest pusty, jest to
Twierdzenie~\ref{thm:chern-gauss-bonnet}.
\end{twierdzenie}
\begin{proof}
Wybierzmy pole wektorowe $V$
równe $\nu$ przy brzegu,
z niezdegenerowanymi zerami $p_i$
we wnętrzu. Przedłużamy
$\nu$ przez kołnierz,
a potem perturbujemy pole
we wnętrzu tak, by było transwersalne
do przekroju zerowego.
Zero transwersalnego pola jest izolowane i niezdegenerowane:
różniczka w punkcie zerowym jest surjekcją między
przestrzeniami tego samego wymiaru, czyli izomorfizmem.
Nie zależy od użytej koneksji, ponieważ wyraz koneksji
mnoży tam $V(p_i)=0$.
Zwartość i brak zer przy brzegu zapewniają, że zer jest skończenie wiele.
Usuńmy małe
rozłączne dyski $D_i$
wokół zer i połóżmy
$u=V/|V|:M_0\to S(TM)$ na pozostałej
części $M_0=M\setminus\bigcup_i\operatorname{int}D_i$.
Ponieważ $\pi\circ u=\id_{M_0}$, mamy
$\dd(u^*\Psi)=u^*(\dd\Psi)=-u^*\pi^*e_{\nabla^{TM}}
=-e_{\nabla^{TM}}|_{M_0}$.
Z Lematu
\ref{lem:angular-form-26}
i twierdzenia Stokesa
otrzymujemy
\[
 \int_{M_0}e_{\nabla^{TM}}
 =-\int_{\partial M_0}u^*\Psi
 =-\int_{\partial M}\nu^*\Psi
   +\sum_i\int_{\partial D_i}u^*\Psi.
\]
Na brzegu dziury znak orientacji
jest przeciwny do orientacji
$\partial D_i$ zorientowanego jako brzeg usuniętego dysku.
Zewnętrzna normalna $M_0$ przy dziurze wskazuje do jej środka,
a normalna zewnętrzna $D_i$ wskazuje od środka.
Te dwa wybory odwracają orientację brzegu.
Gdy dyski
kurczą się do zer, ich
całki zbiegają do stopni
odwzorowań
$u:\partial D_i\to S^{2m-1}$,
czyli indeksów $V$.
Używamy tu definicji
$\operatorname{ind}_{p_i}V=\deg(V/|V|:\partial D_i\to S^{2m-1})$
we współrzędnych dodatniej trywializacji;
niezależność od wyborów omówimy także w
Sekcji~\ref{sec:poincare-hopf-16}.
Aby uzasadnić granicę,
ustalmy zorientowane współrzędne wokół $p_i$
oraz gładką zorientowaną ramkę ortonormalną
$TM$ nad $D_i$ i zapiszmy
$x=p_i+\varepsilon y$, $|y|=1$.
Niech $L$ będzie macierzą $\D_{p_i}V$
w tych dwóch bazach. Współczynniki pola w ramce
spełniają wzór Taylora
$V(p_i+\varepsilon y)=\varepsilon Ly+\varepsilon^2r_\varepsilon(y)$,
gdzie $r_\varepsilon$ i jego pierwsze pochodne po $y$
są jednostajnie ograniczone na sferze.
Odwracalność $L$ daje $c=\min_{|y|=1}|Ly|>0$.
Dla małego $\varepsilon$ mianownik
$|Ly+\varepsilon r_\varepsilon(y)|$ jest więc co najmniej $c/2$.
Stąd
$u_\varepsilon(y)
=(Ly+\varepsilon r_\varepsilon(y))/|Ly+\varepsilon r_\varepsilon(y)|$
zbiega wraz z pierwszymi pochodnymi do $Ly/|Ly|$.
Homotopia zastępująca $\varepsilon$ w liczniku przez
$\tau\varepsilon$, $0\leq\tau\leq1$, także ma niezerowy mianownik.
Stopień $u_\varepsilon$ jest zatem stopniem liniowego modelu,
czyli $\operatorname{sgn}\det L$.
Ostatnia równość wynika z deformacji macierzy w
$\mathrm{GL}^+(2m,\R)$ do identyczności, a w
$\mathrm{GL}^-(2m,\R)$ do jednego odbicia;
na sferze mają one stopnie odpowiednio $1$ i $-1$.
Składowa $\Psi$ zawierająca wyłącznie różniczki
zmiennej włóknowej $u$ jest znormalizowaną
formą objętości sfery, więc jej całka jest
stopniem $u_\varepsilon$. Każdy pozostały
składnik zawiera co najmniej jedną różniczkę
podstawy $\dd x=\varepsilon\dd y$; jego
całka jest $O(\varepsilon)$.
Współczynniki są gładkie i jednostajnie ograniczone,
a pozostałe różniczki $\dd u_\varepsilon$ są ograniczone
na stałej, zwartej sferze parametrów.
Nie pojawia się dodatkowa potęga $\varepsilon^{-1}$:
całkujemy cofniętą formę po tej stałej sferze.
Ponadto całka gładkiej formy Eulera po usuniętych
dyskach dąży do zera jak $O(\varepsilon^{2m})$.
Zatem
\[
 \int_M e_{\nabla^{TM}}
 +\int_{\partial M}\nu^*\Psi
 =\sum_i\operatorname{ind}_{p_i}V.
\]
Suma po prawej wynosi
$\chi(M)$ dla pola
skierowanego na zewnątrz.
Wyprowadźmy używaną
wersję Poincarégo--Hopfa:
wybieramy funkcję Morse'a $f$, dla której
$f=0$ na brzegu, $f<0$ we wnętrzu i $f=-r$
w małym kołnierzu, gdzie $r\geq0$ rośnie do wnętrza.
Jej pochodna w kierunku zewnętrznym jest dodatnia.
Przedłużamy $-r$ do funkcji ujemnej w głębi $M$,
a następnie perturbujemy ją z dala od kołnierza,
zachowując ujemność, aby uzyskać punkty krytyczne Morse'a.
Stosujemy argument z dowodu
Twierdzenia~\ref{thm:gestosc-morse} do zwartego obszaru
poza mniejszym kołnierzem: lokalne funkcje odcinające
wybieramy z nośnikami poza brzegiem.
W samym kołnierzu różniczka $-r$ już jest niezerowa,
więc perturbacja nie jest tam potrzebna.
Wartości krytyczne możemy rozdzielić przez małe zmiany
funkcji w rozłącznych otoczeniach jej skończenie wielu punktów
krytycznych, pozostawiając kołnierz bez zmian.
Jeśli brzeg jest pusty, wybieramy po prostu funkcję Morse'a.
Lokalny argument
dołączenia uchwytu z
Twierdzenia~\ref{thm:uchwyt-morse-17}
stosuje się do kolejnych podpoziomic,
bo w kołnierzu nie ma punktów krytycznych.
Jej gradient jest zewnętrzny, ponieważ
$\langle\operatorname{grad}f,\nu\rangle=\dd f(\nu)>0$.
W punkcie krytycznym różniczka gradientu jest
podniesionym metryką hesjanem:
składniki z pochodną macierzy metryki mnożą $\dd f=0$.
Hesjan ma $\lambda$ ujemnych wartości własnych,
a macierz metryki ma dodatni wyznacznik,
więc indeks pola jest równy $(-1)^\lambda$.
Przyczepienie uchwytu indeksu $\lambda$ dodaje jedną
komórkę wymiaru $\lambda$ i zmienia $\chi$ o ten sam znak.
Przechodząc od pustej podpoziomicy do całego $M$,
otrzymujemy $\sum_p\operatorname{ind}_p(\operatorname{grad}f)=\chi(M)$.

Porównujemy $V$ z $\operatorname{grad}f$ przez
$H_t=(1-t)V+t\operatorname{grad}f$.
Na brzegu składowa normalna tej kombinacji jest dodatnia,
więc żadne zero nie ucieka przez brzeg.
Perturbujemy homotopię
we wnętrzu $M\times(0,1)$,
zachowując oba pola skrajne i brzeg,
do transwersalności
względem przekroju zerowego.
Jest to względne zastosowanie
Wniosku~\ref{cor:transwersalnosc-wzgledna}
do wiązki cofniętej na $M\times[0,1]$.
Zera tworzą wtedy zwartą zorientowaną
$1$-rozmaitość w $M\times[0,1]$,
której końce na poziomie $0$ mają
znaki przeciwne do indeksów pierwszego
pola, a na poziomie $1$ zgodne
z indeksami drugiego pola.
Suma tych znaków na brzegu jest zerem,
więc suma indeksów nie zmienia się
przy homotopii. To kończy dowód.
\end{proof}

\begin{figure}[H]
\centering
\begin{tikzpicture}[>=Stealth]
 \fill[manifoldblue!10] (0,0) circle(2);
 \draw[manifoldblue!65,thick] (0,0) circle(2);
 \foreach \x/\idx in {-.8/1,.8/2}
  {\fill[white] (\x,.15) circle(.44);
   \draw[manifoldgold!85!black,thick] (\x,.15) circle(.44);
   \fill[manifoldblue!80!black] (\x,.15) circle(1.5pt);
   \node[below=2pt,font=\small] at (\x,.15) {$p_{\idx}$};
   \begin{scope}[shift={(\x,.15)}]
    \draw[->,manifoldred,thick] (35:.44)
     arc[start angle=35,end angle=-220,radius=.44];
   \end{scope}}
 \draw[->,manifoldblue,very thick] (20:2)
  arc[start angle=20,end angle=115,radius=2];
 \draw[->,manifoldblue,thick] (145:2)--(145:2.45)
  node[above left=2pt] {$\nu$};
 \node at (0,-1.05) {$M_0$};
 \node[anchor=west,text width=6cm,align=left,font=\small] at (3.1,0)
  {$\partial M_0=\partial M-\sum_i\partial D_i$.\\[6pt]
   Dodatnia orientacja wokół dziury jest przeciwna
   do orientacji brzegu usuniętego dysku.
   To źródło znaku w zastosowaniu Stokesa.};
\end{tikzpicture}
\caption{Orientacje po usunięciu otoczeń zer pola.
Rysunek przedstawia przypadek powierzchni; ta sama reguła
„zewnętrzna normalna jako pierwszy wektor” obowiązuje
w każdym wymiarze parzystym.
Strzałki na okręgach są orientacjami brzegu $M_0$.}
\label{fig:CGB-dziury-orientacja}
\end{figure}

Dla $m=1$ wyprowadźmy związek z krzywizną geodezyjną.
Niech $(e_1,e_2)$ będzie dodatnią ramką ortonormalną,
$Je_1=e_2$, a
$\alpha(X)=g(\nabla_Xe_1,e_2)$ jak w~\eqref{eq:gb-rama}.
Nasza macierz koneksji ma $\omega_{12}=-\alpha$,
czyli $\beta=-\alpha$ w
Przykładzie~\ref{prz:angular-rank-two-17}.
Jeśli $\nu=\cos\theta\,e_1+\sin\theta\,e_2$,
to dodatnia jednostkowa styczna do brzegu jest $T=J\nu$.
Różniczkując w kierunku $T$, dostajemy
$\nabla_T\nu=(\theta'(s)+\alpha(T))T$.
Stąd $\nabla_TT=-(\theta'(s)+\alpha(T))\nu
=(\theta'(s)+\alpha(T))JT$, więc z definicji
$k_g=g(\nabla_TT,JT)=\theta'(s)+\alpha(T)$.
Pochodna $\theta'$ jest liczona względem dodatniego parametru
długości łuku $s$. Wobec tego
$\mathcal B_g=\nu^*\Psi
=(\dd\theta+\alpha)/(2\pi)=k_g\,\dd s/(2\pi)$,
zgodnie z~\eqref{eq:gauss-bonnet-brzeg}.
W ogólnym wymiarze
\eqref{eq:CGB-boundary-form-26}
definiuje formę, która
w~\eqref{eq:chern-gb-brzeg-idea}
była tylko zapowiedziana.
Rozwinięcie transgresji
zawiera drugą formę
fundamentalną brzegu
oraz jego krzywiznę.

\begin{przyklad}[Płaska kula]
\label{prz:CGB-ball-26}
Na $D^{2m}\subset\R^{2m}$
krzywizna znika. Normalna
zewnętrzna ma współczynniki $\nu(u)=u$
na $S^{2m-1}$. W stałej trywializacji płaskiej wiązki
forma kątowa jest znormalizowaną formą objętości
zależną tylko od zmiennej włóknowej.
Po cofnięciu przez normalną dostajemy więc tę samą
formę objętości na sferze brzegowej.
Lemat~\ref{lem:angular-form-26}
daje $\int_{\partial D^{2m}}\mathcal B_g=1
=\chi(D^{2m})$.
Sprawdza to współczynnik
i znak poprawki brzegowej.
W wymiarze dwa ten sam znak sprawdzimy na pierścieniu
$A=\{(x,y):r_0\leq\sqrt{x^2+y^2}\leq R_0\}$,
gdzie $0<r_0<R_0$, z orientacją $\dd x\wedge\dd y$.
Na zewnętrznym okręgu dodatnia orientacja jest przeciwna
do ruchu wskazówek zegara, $k_g=1/R_0$ i $\dd s=R_0\,\dd\theta$.
Na wewnętrznym okręgu orientacja jest zgodna z ruchem
wskazówek zegara.
Parametryzacja $(r_0\cos t,-r_0\sin t)$, $0\leq t\leq2\pi$,
ma jednostkową styczną $T=(-\sin t,-\cos t)$.
Ponieważ $\dd s=r_0\,\dd t$, jej pochodna to
$\dd T/\dd s=(-\cos t,\sin t)/r_0$,
a $JT=(\cos t,-\sin t)$ wskazuje do wnętrza pierścienia,
od środka okręgu.
Zatem $\dd T/\dd s=-(1/r_0)JT$ i $k_g=-1/r_0$.
Wzór daje
\[
 \int_{\partial A}\mathcal B_g
 =\frac1{2\pi}\left(\int_0^{2\pi}\dd\theta
                    -\int_0^{2\pi}\dd t\right)=0=\chi(A).
\]
Ostatnią równość topologiczną daje retrakcja pierścienia
na okrąg: jedna komórka stopnia zero i jedna stopnia jeden.
\end{przyklad}

\begin{przyklad}[Czasza sferyczna: dwa składniki wzoru]
\label{prz:CGB-czasza-17}
Na sferze promienia $R>0$ bierzemy czaszę
$0\leq\rho\leq\rho_0$, gdzie $0<\rho_0\leq\pi/2$
jest kątem od bieguna północnego, a $\varphi\in[0,2\pi]$
jest długością geograficzną.
Metryka $R^2(\dd\rho^2+\sin^2\rho\,\dd\varphi^2)$
z~\eqref{eq:metryka-sfery} daje formę pola
$\dd A=R^2\sin\rho\,\dd\rho\wedge\dd\varphi$.
Z Przykładu~\ref{prz:krzywizna-modele} wiemy, że $K=1/R^2$.
Zatem
\[
 \int_{\text{czasza}}e_{\nabla^{TM}}
 =\frac1{2\pi}\int_0^{2\pi}\int_0^{\rho_0}
          \sin\rho\,\dd\rho\,\dd\varphi
 =1-\cos\rho_0.
\]

\begin{figure}[H]
\centering
\begin{tikzpicture}[>=Stealth]
 \fill[manifoldblue!5] (0,0) circle(1.8);
 \fill[manifoldblue!19] (35:1.8)
  arc[start angle=35,end angle=145,radius=1.8]--cycle;
 \draw[manifoldblue!45] (0,0) circle(1.8);
 \draw[manifoldblue,thick] (35:1.8)
  arc[start angle=35,end angle=145,radius=1.8];
 \draw[manifoldblue!50,dashed] (145:1.8)--(35:1.8);
 \draw[manifoldgold!80!black] (0,0)--(0,1.8);
 \draw[manifoldgold!80!black] (0,0)--(35:1.8);
 \draw[manifoldgold!80!black] (0,.53)
  arc[start angle=90,end angle=35,radius=.53];
 \node at (.42,.67) {$\rho_0$};
 \draw[->,very thick,manifoldred] (35:1.8)--++(-55:.75)
  node[right=3pt] {$\nu$};
 \fill[manifoldblue!80!black] (35:1.8) circle(2pt) (145:1.8) circle(2pt);
 \node[above=4pt] at (0,1.8) {biegun};
 \node[anchor=west,text width=6cm,align=left,font=\small] at (3,0)
  {Wnętrze: $\displaystyle\int_{\text{czasza}}e=1-\cos\rho_0$.\\[6pt]
   Brzeg: $\displaystyle\int_{\partial\text{czasza}}\mathcal B_g=\cos\rho_0$.\\[6pt]
   Zmienia się podział między obie całki,
   a suma pozostaje równa $\chi=1$.};
\end{tikzpicture}
\caption{Czasza sferyczna w przekroju południkowym.
Zaznaczone dwa punkty są przekrojami jej kołowego brzegu;
przerywany odcinek pomaga wskazać jego poziom.
Wektor $\nu$ jest styczny do sfery i wskazuje na zewnątrz czaszy.}
\label{fig:CGB-czasza-dwa-skladniki}
\end{figure}

Na brzegu $\nu=R^{-1}\partial_\rho$,
a dodatnia styczna to $T=(R\sin\rho_0)^{-1}\partial_\varphi$.
Ze wzoru~\eqref{eq:christoffel-metryka} otrzymujemy
$\Gamma^\rho_{\varphi\varphi}
=-\frac12R^{-2}\partial_\rho(R^2\sin^2\rho)
=-\sin\rho\cos\rho$.
Współczynnik pola $T$ jest stały wzdłuż równoleżnika,
więc
$\nabla_TT=-\cos\rho_0/(R^2\sin\rho_0)\,\partial_\rho
=-(\cot\rho_0/R)\nu$.
Z $JT=-\nu$ wynika $k_g=\cot\rho_0/R$.
Ponieważ $\dd s=R\sin\rho_0\,\dd\varphi$, dostajemy
\[
 \int_{\partial\text{czasza}}\mathcal B_g
 =\frac1{2\pi}\int_0^{2\pi}\cos\rho_0\,\dd\varphi
 =\cos\rho_0.
\]
Obie całki sumują się do $1$, czyli charakterystyki
Eulera dysku. Dla półsfery $\rho_0=\pi/2$ brzeg jest
geodezyjnym równikiem i poprawka znika.
Przy $\rho_0\to0$ całka z krzywizny dąży do zera,
a poprawka brzegowa do $1$, jak dla płaskiego dysku.
\end{przyklad}



\par\smallskip
{\footnotesize\noindent\emph{Dalsza lektura:}
\href{https://webhomes.maths.ed.ac.uk/~v1ranick/papers/chern7.pdf}
{S.-S.\ Chern, \emph{A Simple Intrinsic Proof of the
Gauss--Bonnet Formula for Closed Riemannian Manifolds}
(1944), §1--2},
\href{https://www.imar.ro/~sergium/fisiere/pfa.pdf}
{D.\ Cibotaru i S.\ Moroianu,
\emph{Odd Pfaffian Forms},
część o transgresji i brzegu}.\par}
